\chapter{Strategies for the Rest}
\label{c:strategies}

\begin{glance}
The precedent industries teach that the losing move is to defend the rent. Conditional on the book being broadly right, and marking which moves stay correct even if its falsification branches fire: the United States should commoditize before being commoditized; Europe should certify what it cannot manufacture; Japan should build the machine class the war is converging on and referee the trust layer; the Global South should take the commons and unionize the demand side; and China should let its winners charge, consolidate its overcapacity and submit its artifacts to scoreboards it does not control. Every non-principal converges on one institution, the subject of the next chapter.
\end{glance}


\textbf{The United States: win the commoditization race or lose it.} \full{23} US strategy since 2019---chokepoint denial plus closed-frontier rents---produced the outcomes documented here: each denied input spawned a subsidized substitute, and the rent structure concentrated national exposure into the leveraged trade of Chapter~\ref{c:trade}. The counterfactual has never been tried: compete for the commons rather than against it. Fund open-weight frontier models as sustained public infrastructure, not one defensive release; make the American open stack the world's default, because the adoption layer, once lost, prices everything above it; lead rather than restrict RISC-V and open interconnects, which is the Commerce Department's own conclusion~\cite{riscv-restrict}; attack the watt with the urgency now spent on GPU export paperwork; and triage export controls by the model of Section~\ref{c:analogy}. The US capital structure, not its technology, is the strategic vulnerability: an orderly deflation of the scarcity premium, engineered early, is strictly preferable to defending it to the break. That instruction assumed commoditizing and denying were separate policies that could run in parallel; Chapter~\ref{c:mirror} shows they are not, and this is the sharpest finding in the book for American policy. \emph{An export-control regime and an open-weight corporate strategy are aimed at the same object from opposite sides}---American weights released with their training data, an open ISA the same firm has qualified as a CUDA host, and a scale-up domain rebuilt out of merchant Ethernet each lower the cost of exactly the substitution the controls exist to make expensive, and each would have been taken in a world with no Chinese AI industry at all. The conclusion is narrower than that the firms should stop, which would surrender the adoption layer to preserve an instrument whose effectiveness cannot be verified. \emph{A denial policy contingent on the conduct of firms the government does not control is not a strategy}, and the United States must choose which instrument it means: if denial, the controls must reach the conduct that erodes them, which runs into the constitutional difficulty of restraining publication; if commoditization, they should be triaged to the few places they buy real time---EUV, HBM process, advanced packaging---and the energy freed spent on the watt and the open stack. What cannot work is running both at full strength and expecting each to do its job. The counter-evidence bounds it: Hopper shipments to China ran under one per cent of datacentre revenue and the forward outlook assumes none~\cite{nvda-segments} [P], so the cancellation is real and incomplete---the condition under which a choice can still be made deliberately rather than discovered in a filing.

\textbf{Europe: regulation is not production.} Europe enters the fourth war having sat out the previous three's manufacturing and owning the regulatory pen; the trap is repeating solar. Its live assets are ASML---the one true chokepoint, and a wasting one---an open-model tradition, procurement scale, and the credibility to run the certification layer as a neutral. The strategy: build sovereign inference and smaller sovereign-scale training on whichever open stack is best---accept both, audit both, mirror both; spend the AI Act's institutional energy on the trust stack, the one arena where its machinery is an asset; and convert energy policy into compute policy. Certify what you cannot manufacture.

\textbf{Japan: the third-country playbook, first-person.} Japan's assets map one to one onto the book's elements, and five moves follow: build the accelerated-CPU line to its conclusion---an FP8-equipped, capacity-rich successor in the LX2 and Monaka direction, which Japan is one of three polities able to execute [M/S]; adopt \emph{both} commons with the full trust stack, since refusing Chinese weights on principle is strategy-by-flag and taking an American vendor's weights uninspected is the same error wearing the other flag; bid for the certification-authority vacancy, which convened with a small coalition would be the highest-leverage per-yen investment in the entire space; win the performance-per-watt race rather than the watts race; and hedge memory both ways, HBM with allied suppliers and LPDDR-capacity architectures domestically.

\textbf{The Global South: the buyers write the ending.} Take the free stack, keep the exit. Adopt open weights and cheap tokens from whoever supplies them; insist on the trust stack---signed artifacts, local mirrors, evaluation rights---because dependency without auditability is the old colonial bargain in new units; build sovereign inference first and pool regionally for training as capacity-first machines arrive [S]; price national data and deployment access as leverage. Neither financed stack exists in the field (Chapter~\ref{c:world}); the executed template is Brazil's---split the programme, self-finance, train your own people. Whether trust becomes a Chinese utility or a neutral one is the war's last genuinely unclaimed high ground.

\textbf{China: advice for the presumptive winner.} Six failure modes are visible now, most named by Chinese sources. Involution: the model-layer price war is overcapacity without rent recapture; the State Council has warned against ``disorderly competition,'' pricing-law amendments target involution by name, and Xi has asked whether every province needs an AI, compute and EV project~\cite{china-involution}---the fix is at-cost tokens plus paid complements, not the end of competition, which is one of the book's own falsification branches. The monetization gap: post-IPO revenues are thin, Alibaba's adjusted EBITDA fell 84\% in the May 2026 quarter while it committed to exceed a \$56 billion capex plan~\cite{alibaba-capex}, and against a single American lab's \$47 billion run rate the Chinese frontier sector's revenue is a rounding error; let the winners charge, as revenue-threshold licences and metered Max tiers already begin to---rent relocated is the strategy working, rent abolished is the strategy eating itself. Compute overcapacity: up to 80\% of some new capacity unused and rental prices collapsing~\cite{china-idle-dc}---consolidate and publish utilization, because idle capacity is depreciation, not deterrence. The trust deficit: sponsor independent certification of your own artifacts. The physical gaps: HBM, packaging and EUV are absent and CXMT is about three years behind at the class that matters~\cite{china-hbm-gap}. And submit to scoreboards you do not control: no Chinese accelerator has appeared in MLPerf Training, no LX2-class machine has published the AI proof suite, and the flagship's autonomous-silicon-design claim ships without a checkable artifact. Dominating the metrics you control and avoiding the ones you do not is rational for a challenger and corrosive for an incumbent, and China is becoming an incumbent.

